\section{The initial partition of the integer lattice}
\label{sec:area_law_initial_lattice_partition}

% Source: 10-geometry.tex, prop:two-families, lines 154--177,
% 212--218, 299--323 and 545--559,
% revision adc7f1241b42e322a6451854ab7e4b4c146bf78a.
% This is the initial lattice partition before stars and recursive repairs.

\begin{theorem}[Unique initial region at each lattice point]
    \label{thm:area_law_initial_lattice_partition}
    \lean{TNLean.PEPS.AreaLaw.Geometry.exists_unique_initialOpenRegion_integerPoint}
    \leanok
    \uses{def:area_law_initial_open_regions,def:area_law_dyadic_origin}
    Fix the origin $o=(\sqrt2,\sqrt3)$, a nonempty finite endpoint set,
    integer width $C\ge 2$ and initial layer $k_0\ge 50{,}000{,}000$,
    with arbitrary valid residue choices at each layer. Every
    $v\in\mathbb Z^2$ belongs to exactly one actual initial open region
    $\operatorname{int}(X_i^0)$ of
    \cite[Section~11]{OpenAI2026AreaLaw}.
\end{theorem}
\begin{proof}\leanok
    \uses{thm:area_law_initial_closed_cover,
        thm:area_law_initial_lattice_membership,thm:area_law_initial_open_disjoint}
    Closed coverage gives an identifier $i$ with $v\in X_i^0$.
    At a lattice point, birth membership agrees with open membership,
    so $v\in\operatorname{int}(X_i^0)$. If another identifier $j$
    has this property, pairwise disjointness of the open regions forces
    $j=i$.
\end{proof}
